\section{Penelitian Terkait}
\label{sec:related}

\paragraph{Benchmark umum dan benchmark Indonesia} MMLU mengukur pengetahuan multitugas dengan soal pilihan ganda~\cite{hendrycks2021mmlu}; HELM berargumen untuk banyak skenario dan metrik per model, termasuk efisiensi~\cite{liang2023helm}. Untuk bahasa Indonesia, IndoNLU~\cite{wilie2020indonlu} dan IndoLEM~\cite{koto2020indolem} mengumpulkan tugas pemahaman, NusaX~\cite{winata2023nusax} sentimen dan terjemahan untuk bahasa lokal, IndoMMLU~\cite{koto2023indommlu} ujian sekolah---di mana model hanya lulus tingkat sekolah dasar---dan SeaEval~\cite{wang2024seaeval} penalaran multibahasa dan budaya di Asia Tenggara; Cendol melakukan instruction-tuning pada model open-weight untuk bahasa Indonesia~\cite{cahyawijaya2024cendol}. Benchmark-benchmark ini menguji pengetahuan dan pemahaman. SkripsiBench menguji produksi artefak akademik dan kode di baliknya.

\paragraph{Juri LLM} Model bahasa besar yang kuat setuju dengan preferensi manusia terhadap jawaban terbuka~\cite{zheng2023judging,liu2023geval,chiang2023human}, tetapi mereka lebih menyukai generasi mereka sendiri~\cite{panickssery2024selfpref}. Oleh karena itu, kami tidak pernah membiarkan model satu vendor menjadi juri bagi model-model vendor yang sama, menggunakan tiga juri dari tiga vendor yang berbeda, menghilangkan kelunakan masing-masing juri, dan menyiapkan lembar penilaian tersamar untuk membandingkan para juri dengan penilai manusia.

\paragraph{Sitasi dan rujukan} Teks yang dihasilkan dengan sitasi ke kutipan yang diberikan dapat diperiksa secara otomatis~\cite{gao2023alce}, demikian juga ketepatan faktualnya~\cite{min2023factscore}; model bahasa besar diketahui membuat-buat rujukan bibliografis~\cite{walters2023fabrication,huang2025hallucination}. Tugas penulisan berlandaskan sumber kami membatasi sitasi pada enam abstrak yang diberikan dan menandai angka-angka yang tidak ditemukan pada salah satu pun dari abstrak tersebut; tugas rujukan memverifikasi setiap karya yang tercantum dalam basis data bibliografis, sebagaimana yang dilakukan paper-agent untuk manuskripnya sendiri~\cite{arya2026paperagent}.

\paragraph{Kode} HumanEval~\cite{chen2021codex} dan DS-1000~\cite{lai2023ds1000} menguji kode yang dihasilkan dengan unit test. Tugas pengodean kami menjalankan program analisis lengkap pada data sintetis yang menyerupai skripsi dalam sandbox dan membandingkan statistik yang dicetak dengan nilai rujukan dalam batas toleransi.

\paragraph{Biaya} Energi dan biaya model NLP~\cite{strubell2019energy}, dan khususnya inferensi~\cite{luccioni2024power}, berpengaruh pada keputusan penyebaran; kami melaporkan unit penagihan penyedia untuk model API dan energi papan GPU serta waktu yang diukur untuk model lokal, berdampingan dengan kualitas.
